% 并行程序设计 Lab2 — ANN 选题实验报告
% 结构对齐课程《SIMD-LU》雨课堂说明：问题描述 → SIMD 算法设计与复杂性 → 实现
% → 实验与结果分析（含 perf/汇编讨论占位）→ 进阶探索 → Git
% 编译：xelatex 实验报告_Lab2_ANN.tex
\documentclass[UTF8,a4paper,11pt]{ctexart}

\usepackage{geometry}
\geometry{left=2.5cm,right=2.5cm,top=2.5cm,bottom=2.5cm}

\usepackage{amsmath,amssymb}
\usepackage{booktabs}
\usepackage{graphicx}
\usepackage{hyperref}
\usepackage{caption}
\usepackage{enumitem}

\hypersetup{
  colorlinks=true,
  linkcolor=blue,
  urlcolor=blue,
  citecolor=blue
}

\title{并行程序设计实验报告（Lab2 SIMD）\\选题：基于 NEON 的近似最近邻（ANN）检索}
\author{学号：\underline{\hspace{3cm}}\quad 姓名：\underline{\hspace{3cm}}}
\date{\today}

\begin{document}
\maketitle

\begin{abstract}
  本文针对实验指导书 ANN SIMD 选题，在 ARM AArch64（鲲鹏 920，NEON）上实现内积意义下的 Flat、SQ、PQ-ADC、IVF-PQ 及 Top-$R$ 全精度重排等检索路径；给出各阶段 SIMD 算法思路与算术复杂度，汇报 DEEP100K 子集上的延迟与 recall，并预留 perf/汇编分析与进阶对比章节（与《SIMD-LU》评分说明中的基础/进阶要求对应）。
\end{abstract}

% ------------------------------------------------------------------
\section{问题描述}
\label{sec:problem}

\paragraph{任务定义.}
给定基向量集合 $X=\{\mathbf{x}_1,\ldots,\mathbf{x}_N\}\subset\mathbb{R}^d$ 与查询 $\mathbf{q}\in\mathbb{R}^d$，求与 $\mathbf{q}$ 最近的 $k$ 个基向量下标（$k$NN）。本实验采用与指导书一致的内积型伪距离：
\begin{equation}
  \delta(\mathbf{x},\mathbf{y}) = 1 - \sum_{i=1}^{d} x_i y_i ,
\end{equation}
$\delta$ 越小表示内积越大。评测使用 DEEP100K 子集：$N=10^5$，$d=96$；程序对每条查询统计检索耗时并计算 recall@$k$（与 Ground Truth Top-$k$ 交叠比例）。

\paragraph{实验目标（对应 ANN 选题基础要求）.}
在 ARM 平台用 \textbf{NEON Intrinsics} 完成距离/内积热点向量化，结合 OpenMP 等并行策略，对比不同检索结构（Flat、SQ、PQ、IVF-PQ）及参数下的 \textbf{latency--recall}，并用工具从指令与体系结构角度解释性能现象（见第~\ref{sec:exp} 节）。

% ------------------------------------------------------------------
\section{SIMD 算法设计}
\label{sec:design}

\subsection{NEON 浮点内积核（Flat / LUT 共用）}

将 $\sum_{i=1}^d a_i b_i$ 的主循环按宽度 $w\in\{4,8,16\}$ 展开：每次 \texttt{vld1q\_f32} 加载 $a$、$b$ 各一段，用 \texttt{vfmaq\_f32}（或 \texttt{vmlaq\_f32}）累加到 \texttt{float32x4\_t}，最后水平归约；尾部不足 $w$ 的标量收尾。与纯标量相比，理想情况下算术吞吐可提高约 $w$ 倍（受内存带宽与发射端口限制，实测常低于理论值）。

\paragraph{算术复杂度（Flat 暴力扫描）.}
每条查询需对 $N$ 个向量各做一次 $d$ 维内积，串行浮点乘加次数为 $\Theta(Nd)$。SIMD 将内积内层常数因子缩小，但渐近阶不变；全库扫描仍受 $\Theta(Nd)$ 限制。

\subsection{SQ：查询为 float、库为 int8 的混合内积}

量化阶段将每行向量按行尺度压到 int8；查询阶段计算 $\mathrm{scale}_i\cdot \langle \mathbf{q},\hat{\mathbf{x}}_i\rangle$，其中 $\hat{\mathbf{x}}_i$ 为 int8 存储。NEON 将 int8 扩成 float 后与查询向量做 FMA，仍属内积核的变体，复杂度仍为每条查询 $\Theta(Nd)$。

\subsection{PQ-ADC 与 LUT}

将 $\mathbf{x}$ 分为 $m$ 段，每段子空间大小 $d/m$，每段码本大小 $K_s$。对查询 $\mathbf{q}$ 构造 ADC 查找表：
\begin{equation}
  \mathrm{LUT}_{s,c} = \langle \mathbf{q}_s, \mathbf{C}_{s,c}\rangle,\quad
  \hat{\delta}(\mathbf{x},\mathbf{q}) = 1 - \sum_{s=1}^{m} \mathrm{LUT}_{s,\,o_s(\mathbf{x})},
\end{equation}
其中 $o_s(\mathbf{x})$ 为 $\mathbf{x}$ 在第 $s$ 段的 PQ 码字。

\paragraph{复杂度.}
\begin{itemize}[leftmargin=*]
  \item \textbf{LUT 构建}：共 $m K_s$ 个长度为 $d/m$ 的内积，代价 $\Theta(m K_s \cdot (d/m))=\Theta(d K_s)$（与 $m$ 在常用设定下约线性于段数乘子空间维）。
  \item \textbf{全库打分}：对每个库向量累加 $m$ 次 LUT 查表，$\Theta(N m)$；内存访问模式以查表为主，常呈 \textbf{memory-bound}。
\end{itemize}
LUT 填充与 Flat 共用 NEON 内积；可用 OpenMP 对 $(s,c)$ 并行，等价于「对每个子空间并行计算与 $K_s$ 个类中心的内积」。

\subsection{IVF-PQ（结构层）与 SIMD 的分工}

IVF 用全维粗聚类将向量分桶，查询时仅合并 $n_{\mathrm{probe}}$ 个最近桶内 id，再对\textbf{候选子集}做 PQ-ADC。SIMD 仍作用于 PQ 的 LUT 与（若启用）全精度重排内积；粗检索的 query--质心 L$^2$ 本实现主要为标量，属于可进一步向量化的部分。

\subsection{Top-$R$ 全精度重排}

先用 PQ（或 IVF-PQ 候选上的 PQ）得到近似 Top-$R$，再对至多 $R$ 个 id 用与 Flat 相同的 NEON 全维内积重排为 Top-$k$。额外代价约 $\Theta(R d)$，用于换取 recall 大幅提升。

\subsection{实现策略讨论（对应基础要求中的「编程策略」）}

\begin{itemize}[leftmargin=*]
  \item \textbf{对齐}：数据来自 \texttt{fbin}，未强制 16/32 字节对齐时多用 \texttt{vld1q\_f32}（非对齐加载）；若对热点缓冲区做对齐分配，可尝试对齐加载并观察差异（进阶实验可做对比）。
  \item \textbf{向量化粒度}：4/8/16 路封装在寄存器压力与指令数之间折中；$d=96$ 可被 8、16 整除，适合宽向量。
  \item \textbf{并行层次}：OpenMP 用于 LUT 构建与 PQ 全库打分等数据并行段；与 SIMD 互补。
  \item \textbf{Cache}：PQ 打分对 LUT 反复读取、对 \texttt{codes} 顺序扫描，缓存友好性优于随机访问高维向量；IVF 缩小工作集，改善常驻数据局部性。
\end{itemize}

% ------------------------------------------------------------------
\section{实现说明}
\label{sec:impl}

工程为 header-only 模块化设计，主要文件：\texttt{simd\_inner\_product*.h}（NEON 内积）、\texttt{flat\_search\_simd*.h}、\texttt{sq\_search.h}、\texttt{pq\_search.h}（含 ADC LUT、可选 SDC 表）、\texttt{ivf\_pq\_search.h}、\texttt{top\_r\_ip\_rerank.h}。\texttt{main.cc} 中 \texttt{search\_mode} 切换 Flat/SQ/PQ/IVF-PQ 及 Top-$R$ 等模式。编译示例：\texttt{g++ main.cc -O2 -fopenmp -std=c++11 -o main}（平台需开启 NEON）。

% ------------------------------------------------------------------
\section{实验与结果分析}
\label{sec:exp}

\subsection{环境}

实验室 aarch64 服务器（鲲鹏 920），OpenEuler，GCC，OpenMP；数据路径 \texttt{/anndata/}；程序对前 2000 条查询统计平均延迟（$\mu$s）与平均 recall@10。

\subsection{代表性结果}

\begin{table}[htbp]
  \centering
  \caption{IVF-PQ + Top-$R$ 全精度重排（示例：\texttt{./main 5 256 8 1000}）}
  \begin{tabular}{@{}ll@{}}
    \toprule
    指标 & 数值 \\
    \midrule
    索引构建时间 & $\approx 6.42\times 10^{7}\ \mu\mathrm{s}$（约 64.2 s） \\
    平均 Recall@10 & $0.909053$ \\
    平均查询延迟 & $656.171\ \mu\mathrm{s}$ \\
    \bottomrule
  \end{tabular}
\end{table}

该结果体现「粗排（IVF+PQ）缩小候选 + 精排（Top-$R$ IP）修正顺序」的折中：recall 明显高于单阶段 PQ，延迟低于全库 Flat-SIMD。

\subsection{不同算法/参数对比（需补充实测表）}

建议整理小表或曲线：\textbf{Flat（mode 0/3/4）}、\textbf{SQ}、\textbf{PQ（mode 2）}、\textbf{IVF-PQ（mode 5）} 及不同 $R$、$n_{\mathrm{probe}}$ 的 recall--latency；说明共享机器上宜多次运行取中位数。

\subsection{性能剖析与汇编（基础要求建议完成项）}

\paragraph{perf.}
对 Flat 与 PQ 路径分别采样，关注 \texttt{cycles}、\texttt{instructions}、\texttt{cache-misses}、\texttt{branch-misses}。PQ 阶段若 IPC 较高但绝对时间长，多因 $N$ 大且访存主导；Flat 若 L1/L2 miss 高，与全维向量流式读取有关。

\paragraph{编译器汇编.}
使用 \texttt{-S -fverbose-asm -O2} 或 \texttt{objdump -d} 对比热点循环是否生成 \texttt{fmla}（NEON FMA）与展开程度；说明手写 intrinsics 与编译器对纯标量循环自动向量化的差异（可与第~\ref{sec:advanced} 节进阶内容合并）。

% ------------------------------------------------------------------
\section{进阶要求（选做，对应《SIMD-LU》第三节）}
\label{sec:advanced}

\begin{enumerate}[leftmargin=*]
  \item \textbf{手写 SIMD 与自动向量化对比}：对同一标量内积循环分别采用 \texttt{-O2/-O3} 不加 intrinsics 与手写 NEON 版本，对比汇编中向量指令比例与实测耗时，\textbf{分析差异来源}（不仅罗列汇编）。
  \item \textbf{跨平台对比}：若有 x86 环境，对比 SSE/AVX 与 ARM NEON 在同一代码逻辑下的性能及架构因素（向量宽度、缓存、频率）；若仅有 ARM，可说明限制并对比不同编译选项。
  \item \textbf{其他优化}：如 IVF 粗 L2 的 NEON、PQ Fast Scan 布局、SDC/两阶段重排等（结合个人实现情况撰写）。
  \item \textbf{生成式 AI 辅助学习}：将符合课程要求的对话记录（问题—思考—改进，而非直接索要完整答案）整理为附录。
\end{enumerate}

% ------------------------------------------------------------------
\section{总结}

本报告按课程 SIMD 实验评分说明组织：完成了 ANN 选题下的 NEON 内积核、SQ/PQ/IVF-PQ 及 Top-$R$ 重排等设计与实现，给出主要步骤的复杂性讨论，并汇报 IVF-PQ 重排下的代表性 recall 与延迟。后续可在第~\ref{sec:exp}、\ref{sec:advanced} 节补全 perf/汇编截图、参数扫描表及进阶对比，以符合「深入分析」的加分标准。

\vspace{1em}
\noindent\textbf{Git 仓库：}\url{https://example.com/your-repo}（请替换为实际链接）

\appendix
\section{附录：截图与补充材料}

可将终端运行截图、\texttt{perf report} 截图、关键汇编片段等置于本附录。插图示例：
\begin{verbatim}
\includegraphics[width=0.92\linewidth]{screenshot_ivf_pq.png}
\end{verbatim}

\end{document}
